% (This file is included by thesis.tex; you do not latex it by itself.)

\begin{abstract}

% The text of the abstract goes here.  If you need to use a \section
% command you will need to use \section*, \subsection*, etc. so that
% you don't get any numbering.  You probably won't be using any of
% these commands in the abstract anyway.

\noindent This thesis develops probabilistic methods for studying how discrete stochastic systems can be
sampled, inferred, and approximated. Across problems in phylogenetics, adaptive sampling, combinatorial
growth, and stochastic control, we exploit the structure of the underlying stochastic processes to obtain
finite-sample guarantees, characterize optimal procedures, and identify informative asymptotic limits.

% Bipartitions
\noindent In Chapter 1, we study bipartition coverage probabilities under the multispecies coalescent model. 
These bipartitions are used by some species-tree inference procedures to restrict the relevant search space,
and thus our bounds provide meaningful estimates of the sample complexity of these restricted algorithms. 
Practically, our bounds remain below biologically realistic numbers of loci across a substantially broader range of parameter 
settings, expanding their usefulness for empirical datasets. Theoretically, our analysis sharpens
our understanding of coalescence under the MSC model and develops new asymptotics for these bounds
and absorption times under Kingman's coalescent in the natural short branch regime.

% Inference under the PCH model
\noindent In Chapter 2, we study phylogenetic inference under the 
PCH model of \cite{canby2024_polymorphism_linguistic_phylogenetics}, a stochastic model for the
evolution of linguistic characters which explicitly incorporates polymorphism. 
The foundation of this model is a simple birth-death process which describes how 
states evolve down the underlying phylogenetic tree. Building off the work of \cite{evans2026_canby_identifiability}, 
we analyze various probabilistic properties of this process, translating them into measures of quartet support
and candidate bipartitions relevant to phylogenetic inference. These results motivate ASTRAL-based 
inference procedures and provide preliminary sample-complexity bounds for recovering 
the underlying language phylogeny from independently evolving characters.

% Self-Balancing Sampling
\noindent In Chapter 3, in joint work with Daniel Raban we study a family of sequential sampling rules 
that adaptively bias sampling probabilities in order to achieve faster convergence of the empirical distribution 
to a desired target law, while keeping the resulting samples as unpredictable as possible. The resulting scheme 
arises naturally among a class of Markovian samplers sharing a certain translation-invariance, converges at the fastest
possible $O(n^{-1})$ rate, and is the unique solution to a natural entropy-regularized optimization problem.
We characterize the influence of a natural bias parameter on the trade-off between convergence and unpredictability,
as measured by an adversial prediction accuracy.

% Exchangeable arrays
In Chapter 4, we connect the modern theory of hypergraphons to the classical theory of exchangeable arrays
and Doob-Martin boundary theory. We show a broad range of combinatorial growth processes can be encoded in
such arrays, and that the graphon notion of left-convergence via embedding densities easily generalizes. Moreover,
we show the generated topology is equivalent to the Doob-Martin topology, generalizing some of the work of
\cite{grubel2015_persisting_randomness_graphs_search_trees,grubel2023_ranks_copulas_permutons}. 

% Diffusion Limits Entropy-Regularized Control
Lastly, in Chapter 5, we provide a more natural derivation of the continuous-time limit to an 
entropy-regularized optimal control problem, initially proposed in \cite{wang2020_reinforcement_learning_stochastic_control}.
Our construction provides path-wise convergence of controlled trajectories and local uniform convergence
of the associated value functions. 

\end{abstract}
